\begin{table}[H]
\small
\setlength{\tabcolsep}{4pt}

\caption{\label{tab:ea20-policy-method-sensitivity}Sensitivity of euro-area inflation-inequality effects to counterfactual construction methods}
\centering
\resizebox{\textwidth}{!}{%
\begin{tabular}[t]{lrrrrrr}
\toprule
\multicolumn{1}{c}{} & \multicolumn{6}{c}{Inflation inequality}\\
\cmidrule(lr){2-7}
\multicolumn{1}{c}{} & \multicolumn{3}{c}{Individual P20--P80} & \multicolumn{3}{c}{Income Q1--Q5}\\
\cmidrule(lr){2-4} \cmidrule(lr){5-7}
Scenario & 2021 & 2022 & 2023 & 2021 & 2022 & 2023\\
\midrule
Full specification & -0.1 & -1.8 & -0.7 & -0.1 & -0.6 & -0.6\\
Direct reconstructions only & 0.0 & -1.1 & -0.2 & 0.0 & -0.4 & -0.1\\
Full specification without conditional benchmarks & -0.1 & -1.7 & -0.4 & -0.1 & -0.6 & -0.3\\
Conditional benchmark wedges at 50\% & -0.1 & -1.7 & -0.5 & -0.1 & -0.6 & -0.4\\
Conditional benchmark wedges at 150\% & -0.1 & -1.8 & -0.9 & -0.1 & -0.7 & -0.7\\
\addlinespace
Dependence on estimated layers: full minus direct-only & -0.1 & -0.7 & -0.4 & 0.0 & -0.3 & -0.4\\
\bottomrule
\end{tabular}
}%
\begin{flushleft}
\footnotesize{Notes. The first block reports the observed minus counterfactual P20--P80 range of annual household inflation, in percentage points. The range is P80 minus P20. National percentile endpoints use HBS survey weights and are averaged with annual HICP country weights over the 19-country household sample; Portugal is excluded. The second block reports the corresponding Q1--Q5 effects for all 20 euro-area countries. Direct reconstructions use statutory rates, unit rebates, or published reference-price paths without a fitted intensity. The 50\% and 150\% rows scale only the conditional benchmark price wedges. Full minus direct-only measures dependence on estimated layers; it is not a confidence interval. Sources: Eurostat, HBS 2020 microdata, national sources, authors' calculations.}
\end{flushleft}
\end{table}
